\documentclass[11pt,a4paper]{article}

\usepackage[T1]{fontenc}
\usepackage[utf8]{inputenc}
\usepackage{lmodern}
\usepackage[margin=2.5cm]{geometry}
\usepackage{amsmath,amssymb,amsthm,mathtools}
\usepackage{bm}
\usepackage{booktabs}
\usepackage{array}
\usepackage{xcolor}
\usepackage{listings}
\usepackage{enumitem}
\usepackage[colorlinks=true,linkcolor=blue!50!black,citecolor=blue!50!black,urlcolor=blue!50!black]{hyperref}

\lstdefinestyle{cpp}{
  language=C++,
  basicstyle=\ttfamily\footnotesize,
  keywordstyle=\color{blue!60!black}\bfseries,
  commentstyle=\color{green!40!black}\itshape,
  stringstyle=\color{red!50!black},
  showstringspaces=false,
  breaklines=true,
  frame=single,
  framerule=0.3pt,
  rulecolor=\color{black!30},
  backgroundcolor=\color{black!3},
  xleftmargin=1em,
  framexleftmargin=0.5em,
  columns=fullflexible,
  keepspaces=true
}
\lstset{style=cpp}
\setlength{\emergencystretch}{3em}

\newtheorem{proposition}{Proposition}
\theoremstyle{remark}
\newtheorem*{remark}{Remark}

\newcommand{\R}{\mathbb{R}}
\newcommand{\Z}{\mathbb{Z}}
\newcommand{\T}{^{\mathsf T}}
\newcommand{\code}[1]{\texttt{#1}}

\title{Discrete Burgers Circuit Dislocation Analysis (Method 2)\\[4pt]
\large Theoretical Formulation and Implementation in the MTM Lattice-Triangulation Code}
\author{Technical Note}
\date{October 2026}

\begin{document}
\maketitle

\begin{abstract}
In 2D nonlinear finite-element and atomistic crystal plasticity simulations on Delaunay
triangulations, topological vertex valence (\code{NodeValence} / \code{ReferenceCoordination},
counting triangles incident to each atom) works on six-fold triangular lattices ($p=6$)
where every elementary triangle is equilateral, but breaks down on four-fold square lattices
($p=4$). Because each $a_0 \times a_0$ square unit cell is split by a diagonal of length
$\sqrt{2}\,a_0$ chosen between two degenerate orientations (\code{/} and \code{\textbackslash}),
any boundary between \code{/} and \code{\textbackslash} diagonal domains produces vertices of
triangle valence $5$ and $7$ in a $100\%$ defect-free crystal. This note presents the
theoretical formulation and implementation of \textbf{Method 2 (Discrete Burgers Circuit)},
a discrete exterior calculus formulation of Nye's dislocation density tensor
$\bm{\alpha} = -\operatorname{Curl}(\bm{F}_e^{-1})$. By projecting directed mesh edges onto
the ideal Bravais lattice as an anti-symmetric $\Z^2$-valued $1$-cochain $\Delta\bm{n}_{ab}$
and taking its discrete exterior derivative $\bm{B}_T = (d\Delta\bm{n})_T$ around each
counter-clockwise triangle $T$, Method 2 guarantees $\bm{B}_T \equiv \bm{0}$ on every
defect-free crystal configuration (regardless of \code{/} vs.\ \code{\textbackslash} diagonal
switches or integer surface steps), satisfies Discrete Stokes' Theorem identically across arbitrary
patches and dislocation junctions, and localizes dislocation cores to single elementary triangles.
\end{abstract}

\tableofcontents

\section{Why Topological Triangle Valence Fails on Square Lattices}

Let $\mathcal{T} = (\mathcal{V}, \mathcal{E}, \mathcal{F})$ be a Delaunay triangulation of
$N$ atoms with deformed positions $\{\bm{y}_i\}_{i=1}^N \subset \R^2$. The topological
triangle coordination number of vertex $i \in \mathcal{V}$ is the number of incident triangles:
\begin{equation}
  z_i^{\mathrm{tri}} = \#\{T \in \mathcal{F} : i \in T\}.
\end{equation}
In a 2D triangular (hexagonal) Bravais lattice ($p=6$) with lattice spacing $a_0$, each bulk
atom has $6$ nearest neighbors at distance $a_0$, every elementary triangle is equilateral
with side lengths $(a_0, a_0, a_0)$, and $z_i^{\mathrm{tri}} = 6$ everywhere in a defect-free
crystal. An isolated edge dislocation corresponds to a $5/7$ coordination dipole ($z_i^{\mathrm{tri}}=5$
adjacent to $z_j^{\mathrm{tri}}=7$).

In a 2D square Bravais lattice $\Z^2 a_0$ ($p=4$), however, each atom has only \textbf{four}
nearest neighbors at distance $a_0$ and \textbf{four} second-nearest neighbors at distance
$\sqrt{2}\,a_0$. Every elementary $a_0 \times a_0$ square plaquette is cocircular (its four
vertices lie on a circle of diameter $\sqrt{2}\,a_0$) and must be split into two right-angled
isosceles triangles of side lengths $(a_0, a_0, \sqrt{2}\,a_0)$ by inserting either:
\begin{itemize}[noitemsep]
  \item the \textbf{positive diagonal} (\code{/}) connecting $(m,n) \leftrightarrow (m+1,n+1)$, or
  \item the \textbf{negative diagonal} (\code{\textbackslash}) connecting $(m+1,n) \leftrightarrow (m,n+1)$.
\end{itemize}
At any interior atom $i$, the four surrounding square cells always contribute $4$ axis-aligned
nearest-neighbor bonds, plus $d_i \in \{0, 1, 2, 3, 4\}$ diagonal bonds of length $\sqrt{2}\,a_0$
that happen to touch atom $i$. Hence the triangle valence at atom $i$ is:
\begin{equation}
  z_i^{\mathrm{tri}} = 4 + d_i \in \{4, 5, 6, 7, 8\}.
\end{equation}
Whenever floating-point roundoff or minuscule elastic strain causes adjacent square cells to
switch between \code{/} and \code{\textbackslash} diagonals, $d_i$ takes values $1$ and $3$,
producing $z_i^{\mathrm{tri}} = 5$ and $7$ in a completely defect-free square lattice.

\section{Continuum vs.\ Discrete Multiplicative Plasticity (\texorpdfstring{$\bm{F} = \bm{F}_e \bm{F}_p$}{F = Fe Fp})}

\subsection{Kr\"oner--Bilby Decomposition and Nye's Dislocation Tensor}

In the Mesoscopic Tensorial Model (MTM) on a fixed reference triangulation, the total element
deformation gradient $\bm{F} = \nabla\bm{y}$ is decomposed multiplicatively as
\begin{equation}
  \bm{F} = \bm{F}_e \bm{F}_p,
\end{equation}
where $\bm{F}_p \in \mathrm{GL}(2,\Z)$ is the lattice-invariant plastic slip obtained via
Lagrange reduction $\bm{C}_{\mathrm{red}} = \bm{M}\T \bm{C} \bm{M}$ with $\bm{M} = \bm{F}_p^{-1} \in \mathrm{SL}(2,\Z)$
($\bm{C} = \bm{F}\T \bm{F}$), and $\bm{F}_e = \bm{F}\bm{F}_p^{-1} = \bm{F}\bm{M}$ is the
\textbf{elastic distortion tensor}.

Pulling back an infinitesimal deformed line element $d\bm{y}$ by $\bm{F}_e^{-1}$ yields the
local elastic lattice differential:
\begin{equation}
  d\bm{\xi} = \bm{F}_e^{-1}\,d\bm{y}.
\end{equation}
In a compatible (defect-free) crystal, $d\bm{\xi}$ is an exact differential ($\operatorname{Curl}(\bm{F}_e^{-1}) = \bm{0}$).
Around a dislocation core, $d\bm{\xi}$ is no longer closed, and its circulation around a
counter-clockwise contour $\mathcal{C} = \partial\Omega$ defines the true continuum Burgers vector:
\begin{equation}
  \bm{b}(\mathcal{C}) = -\oint_{\mathcal{C}} \bm{F}_e^{-1}\,d\bm{y}
  = -\int_{\Omega} \operatorname{Curl}(\bm{F}_e^{-1})\,dA,
  \label{eq:continuum_burgers}
\end{equation}
where $\bm{\alpha} = -\operatorname{Curl}(\bm{F}_e^{-1})$ is \textbf{Nye's dislocation density tensor}.

\subsection{Effect of Topological Remeshing on \texorpdfstring{$\bm{F}_p$}{Fp}}

When adaptive Delaunay remeshing (\code{AdaptiveMesher}) is enabled in the MTM code, topological
edge flips reconnect atoms to their \emph{new} physical neighbors after each plastic slip event
 and reset the element reference configuration. Consequently, in a region where a dislocation has
already glided through and left a slipped perfect crystal behind, the remeshed element deformation
gradients satisfy $\bm{F}_T \approx \bm{I}$ ($\bm{M}_T = \bm{I}$): the plastic slip $\bm{F}_p$ has
been absorbed into the \textbf{graph connectivity} of the mesh.

Method 2 evaluates the discrete counterpart of Eq.~\eqref{eq:continuum_burgers} directly on the
directed edges of the current triangulation $\mathcal{T}$, making it invariant to whether plastic
slip was accommodated by $\bm{M}_T \in \mathrm{SL}(2,\Z)$ or by Delaunay reconnection.

\section{Discrete Exterior Calculus Formulation of Method 2}

\subsection{Reference Frame Pull-Back and Automatic Bravais Basis \texorpdfstring{$\bm{A}_0$}{A0}}

For each vertex $a$ of an active triangle $T \in \mathcal{F}$ with node index $n_a$ and periodic
translation vector $\bm{t}_a$, we first pull back the position by the macroscopic external
deformation gradient $\bm{F}_{\mathrm{ext}}$:
\begin{equation}
  \bm{x}_a = \bm{F}_{\mathrm{ext}}^{-1}\,\bm{y}_{n_a} + \bm{t}_a \in \R^2.
\end{equation}
Each triangle $T = (0, 1, 2)$ is oriented counter-clockwise in the reference frame so that
$(\bm{x}_1 - \bm{x}_0) \times (\bm{x}_2 - \bm{x}_0) > 0$.

The ideal Bravais lattice basis matrix $\bm{A}_0 \in \R^{2\times 2}$ of the crystal is defined as
\begin{equation}
  \bm{A}_0 = a_0\,\bm{R}(\theta_0)\,\bm{B}_{\mathrm{lat}}, \qquad
  \bm{R}(\theta_0) = \begin{pmatrix} \cos\theta_0 & -\sin\theta_0 \\ \sin\theta_0 & \cos\theta_0 \end{pmatrix},
\end{equation}
where the primitive dimensionless basis $\bm{B}_{\mathrm{lat}}$ is
\begin{equation}
  \bm{B}_{\mathrm{lat}}^{\mathrm{square}} = \begin{pmatrix} 1 & 0 \\ 0 & 1 \end{pmatrix} \quad (p=4),
  \qquad
  \bm{B}_{\mathrm{lat}}^{\mathrm{triangular}} = \begin{pmatrix} 1 & 1/2 \\ 0 & \sqrt{3}/2 \end{pmatrix} \quad (p=6).
\end{equation}
The parameters $(p, a_0, \theta_0)$ are extracted automatically from $\mathcal{T}$:
\begin{enumerate}[noitemsep]
  \item For each triangle $T$ with sorted edge lengths $L_1(T) \le L_2(T) \le L_3(T)$, the median
        ratio $\operatorname{median}_T(L_3/L_1)$ separates triangular lattices ($\approx 1.0 < 1.20$,
        $p=6$) from square lattices ($\approx \sqrt{2} \approx 1.414 > 1.20$, $p=4$).
  \item The lattice constant $a_0$ is taken from \code{userData->ideal\_lattice\_parameter} (with
        fallback to $\operatorname{median}_T(L_1)$).
  \item The global crystal orientation angle $\theta_0 \in (-\pi/p, \pi/p]$ is computed from all
        nearest-neighbor edges ($L_{ab} \in [0.75\,a_0, 1.22\,a_0]$) via the $p$-fold bond-orientational
        order parameter:
        \begin{equation}
          Z_p = \sum_{(a,b) \in \mathcal{E}_{\mathrm{NN}}} e^{i p \phi_{ab}}, \qquad
          \theta_0 = \frac{1}{p}\operatorname{atan2}(\operatorname{Im} Z_p, \operatorname{Re} Z_p).
        \end{equation}
\end{enumerate}

\subsection{Canonical Anti-Symmetric Lattice Step \texorpdfstring{$1$}{1}-Cochain \texorpdfstring{$\Delta\bm{n}_{ab}$}{Delta n\_ab}}

On every directed edge $e_{ab} = (a \to b) \in \mathcal{E}$, we map the relative vector
$\bm{x}_b - \bm{x}_a$ into the ideal Bravais basis and round to the nearest integer lattice vector
$\Delta\bm{n}_{ab} \in \Z^2$. To ensure that $\Delta\bm{n}$ is a strict \textbf{algebraic $1$-cochain}
($\Delta\bm{n}_{ba} \equiv -\Delta\bm{n}_{ab}$ even when a component lies near a half-integer tie
$\pm 0.5$), we canonicalize the edge direction using the lexicographical order
$(n_a, x_a, y_a) \prec (n_b, x_b, y_b)$:
\begin{equation}
  \Delta\bm{n}_{ab} =
  \begin{cases}
    \operatorname{round}\!\left(\bm{A}_0^{-1}(\bm{x}_b - \bm{x}_a)\right) & \text{if } a \prec b, \\[4pt]
    -\operatorname{round}\!\left(\bm{A}_0^{-1}(\bm{x}_a - \bm{x}_b)\right) & \text{if } b \prec a.
  \end{cases}
  \label{eq:canonical_step}
\end{equation}

\subsection{Elementary Triangle Burgers Circuit (\texorpdfstring{$2$}{2}-Cochain \texorpdfstring{$\bm{B}_T$}{B\_T})}

For each counter-clockwise triangle $T = (0, 1, 2) \in \mathcal{F}$ (excluding non-physical outer
boundary slivers with $L_3(T) > 2.15\,a_0$), the discrete exterior derivative $d\Delta\bm{n}$
around $\partial T$ defines the integer Burgers vector of triangle $T$:
\begin{equation}
  \bm{B}_T = (d\Delta\bm{n})(T) = \Delta\bm{n}_{01} + \Delta\bm{n}_{12} + \Delta\bm{n}_{20} \in \Z^2,
  \label{eq:triangle_burgers}
\end{equation}
and the dimensionless physical Burgers vector (in units of $a_0$):
\begin{equation}
  \bm{b}_T = \frac{1}{a_0}\bm{A}_0 \bm{B}_T \in \R^2, \qquad b_T = \|\bm{b}_T\|_2.
\end{equation}

\begin{proposition}[Exact Vanishing on Defect-Free Square and Triangular Lattices]
Let $\Omega \subset \mathcal{T}$ be any defect-free crystal region in which each atom $a \in \Omega$
is displaced from an integer Bravais site $\bm{m}_a \in \Z^2$ by an elastic displacement $\bm{u}_a$
($\bm{x}_a = \bm{A}_0 \bm{m}_a + \bm{u}_a$) satisfying the elastic compatibility bound
$\|\bm{A}_0^{-1}(\bm{u}_b - \bm{u}_a)\|_\infty < \tfrac{1}{2}$ on every edge $(a,b)$ in $\Omega$.
Then $\bm{B}_T \equiv (0,0)\T$ for every triangle $T \subset \Omega$, regardless of whether square
diagonals are oriented as \code{/} or \code{\textbackslash} or whether boundaries contain integer
surface steps.
\end{proposition}
\begin{proof}
Since $\|\bm{A}_0^{-1}(\bm{u}_b - \bm{u}_a)\|_\infty < \tfrac{1}{2}$, Eq.~\eqref{eq:canonical_step}
evaluates to the exact lattice difference $\Delta\bm{n}_{ab} = \bm{m}_b - \bm{m}_a = (d\bm{m})_{ab}$.
Because $\Delta\bm{n} = d\bm{m}$ is an exact $1$-coboundary and $d^2 \equiv 0$, summing around
any triangle $T = (0,1,2)$ yields
\[
  \bm{B}_T = (\bm{m}_1 - \bm{m}_0) + (\bm{m}_2 - \bm{m}_1) + (\bm{m}_0 - \bm{m}_2) \equiv \bm{0}. \qedhere
\]
\end{proof}

\subsection{Discrete Stokes' Theorem, Plaquette Cancellation, and Junctions}

\begin{proposition}[Discrete Stokes' Theorem and Frank's Rule]
For any patch of triangles $\Omega = \bigcup_{k=1}^M T_k \subset \mathcal{F}$ with counter-clockwise
outer boundary loop $\partial\Omega$, the circulation of $\Delta\bm{n}$ around $\partial\Omega$
equals the exact integer sum of elementary triangle Burgers vectors inside $\Omega$:
\begin{equation}
  \bm{B}(\partial\Omega) \equiv \oint_{\partial\Omega} \Delta\bm{n}_{ab}
  = \sum_{T \subset \Omega} \bm{B}_T \in \Z^2.
  \label{eq:discrete_stokes}
\end{equation}
\end{proposition}
\begin{proof}
Every interior edge shared by two triangles $T_1, T_2 \subset \Omega$ is traversed once as $a \to b$
and once as $b \to a$. By Eq.~\eqref{eq:canonical_step}, $\Delta\bm{n}_{ab} + \Delta\bm{n}_{ba} \equiv \bm{0}$
in exact integer arithmetic, leaving only the outer boundary edges $\partial\Omega$.
\end{proof}

Two key consequences follow directly from Proposition~2:
\begin{enumerate}
  \item \textbf{Unit-Cell Plaquette Cancellation Across Shared Square Diagonals:}
        In a square lattice, two triangles $T_1, T_2$ sharing a diagonal edge of length
        $L_{\mathrm{shared}} \approx \sqrt{2}\,a_0 > 1.15\,a_0$ form a single $4$-atom unit-cell
        plaquette $Q = T_1 \cup T_2$ containing no interior atoms. If only the shared internal
        diagonal edge crosses a rounding threshold under elastic shear while the four outer
        unit-cell bonds remain compatible, then $\bm{B}(Q) = \bm{B}_{T_1} + \bm{B}_{T_2} = \bm{0}$
        with $\bm{B}_{T_1} = -\bm{B}_{T_2} \neq \bm{0}$. Cancelling such opposite pairs
        ($\bm{B}_{T_1} = \bm{B}_{T_2} = \bm{0}$) removes single-diagonal elastic rounding artifacts
        while leaving $\sum_{T \subset \Omega} \bm{B}_T$ invariant on every larger circuit.
  \item \textbf{Nodal $1$-Ring Circuits and Dislocation Junctions:}
        Summing $\bm{B}_T$ over the $1$-ring $\Omega_i = \bigcup_{T \ni i} T$ of each atom $i$ yields
        the nodal $1$-ring Burgers vector:
        \begin{equation}
          \bm{B}_i^{(1\text{-ring})} = \sum_{T \ni i} \bm{B}_T \in \Z^2, \qquad
          \bm{b}_i^{(1\text{-ring})} = \frac{1}{a_0}\bm{A}_0 \bm{B}_i^{(1\text{-ring})}.
        \end{equation}
        When two dislocations with Burgers vectors $\bm{B}_1$ and $\bm{B}_2$ intersect or react to
        form a junction (e.g.\ $\bm{B}_1 = (1,0)\T$ and $\bm{B}_2 = (0,1)\T$ in \code{Shift / Staircase}),
        Eq.~\eqref{eq:discrete_stokes} automatically enforces Frank's conservation rule
        $\bm{B}_{\mathrm{junction}} = \bm{B}_1 + \bm{B}_2 = (1,1)\T$.
\end{enumerate}

\section{Summary of Exported ParaView VTK Fields}

Table~\ref{tab:vtk_fields} lists the VTK fields written by \code{ConfigurationSaver::writeToVTKFile}
and \code{ConfigurationSaver::writeToVTK\_DefectAnalysis}.

\begin{table}[htbp]
  \centering
  \caption{Method 2 Discrete Burgers Circuit fields exported to binary VTK files.}
  \label{tab:vtk_fields}
  \vspace{4pt}
  \begin{tabular}{@{}lllp{7.2cm}@{}}
    \toprule
    \textbf{VTK Field Name} & \textbf{Section} & \textbf{Type} & \textbf{Physical Meaning} \\
    \midrule
    \code{ElementBurgersMagnitude} & \code{CELL\_DATA} & \code{float} & $\|\bm{b}_T\|_2$ per triangle in units of $a_0$ ($0$ in perfect crystal, $1$ at elementary dislocation core) \\
    \code{ElementBurgersVector} & \code{CELL\_DATA} & \code{float[3]} & $(b_{T,x}, b_{T,y}, 0)$ per triangle in units of $a_0$ \\
    \code{BurgersMagnitude} & \code{POINT\_DATA} & \code{float} & $\|\bm{b}_i^{(1\text{-ring})}\|_2$ around the $1$-ring of atom $i$ \\
    \code{BurgersVector} & \code{POINT\_DATA} & \code{float[3]} & $(b_{i,x}^{(1\text{-ring})}, b_{i,y}^{(1\text{-ring})}, 0)$ around the $1$-ring of atom $i$ \\
    \code{ReferenceCoordination} & \code{POINT\_DATA} & \code{float} & Burgers-filtered coordination ($6$ whenever $\bm{B}_i^{(1\text{-ring})}=\bm{0}$; $5/7$ only at $\bm{B}_i^{(1\text{-ring})}\neq\bm{0}$ cores) \\
    \code{NodeValence} & \code{POINT\_DATA} & \code{float} & Unfiltered mesh triangle valence $z_i^{\mathrm{tri}}$ \\
    \bottomrule
  \end{tabular}
\end{table}

\end{document}
